\section{Bound With Orientation Uncertainty}
\label{app:orientation}
Let $\vect n_i(\bm\eta)$ be the true normals, where the unknown vector $\bm\eta$ describes their small rotations away from nominal normals $\vect n_i=\vect n_i(\vect 0)$. Two sets of observations are available: the RSS, $\vect y=\vect s(\vect r,\bm\eta)+\bm\epsilon$ with $\bm\epsilon\sim\mathcal N(\vect 0,\vect W^{-1})$, and the sensor reading $\hat{\bm\eta}=\bm\eta+\vect e$ with $\vect e\sim\mathcal N(\vect 0,\bm\Sigma_\theta)$ independent of $\bm\epsilon$; the reported normals are $\vect n_i(\hat{\bm\eta})$. Up to a constant, the joint log-likelihood is $-\frac12(\vect y-\vect s)\trans\vect W(\vect y-\vect s)-\frac12(\hat{\bm\eta}-\bm\eta)\trans\bm\Sigma_\theta^{-1}(\hat{\bm\eta}-\bm\eta)$ and, with $\vect J=\partial\vect s/\partial\vect r\trans$ and $\vect D=\partial\vect s/\partial\bm\eta\trans$, the FIM of $(\vect r,\bm\eta)$ is~\cite{Kay1993}:
\begin{equation}
 {\vect I_{\vect r,\bm\eta}=\begin{bmatrix}
  \vect J\trans\vect W\vect J & \vect J\trans\vect W\vect D\\
  \vect D\trans\vect W\vect J & \vect D\trans\vect W\vect D+\bm\Sigma_\theta^{-1}
 \end{bmatrix}.}
 \label{eq:joint-fim}
\end{equation}
The bound on $\vect r$ is the upper-left $3\times3$ block of the inverse of~\eqref{eq:joint-fim}, that is, the inverse of the Schur complement of its lower-right block:
\begin{equation}
 {\begin{aligned}
  \vect I_e&=\vect J\trans\vect W\vect J-\vect J\trans\vect W\vect D\left(\vect D\trans\vect W\vect D+\bm\Sigma_\theta^{-1}\right)^{-1}\vect D\trans\vect W\vect J\\
  &=\vect J\trans\left(\vect W^{-1}+\vect D\bm\Sigma_\theta\vect D\trans\right)^{-1}\vect J,
 \end{aligned}}
 \label{eq:schur}
\end{equation}
where the second line follows from the matrix inversion lemma. This is~\eqref{eq:fim-orientation}. It is evaluated at the reported normals, which differ from the true ones only by the small errors, and it bounds any unbiased estimator of $\vect r$ that uses the RSS and the sensor readings. The estimators of Section~\ref{subsec:estimator} take the reported normals as exact and need not attain it (Section~\ref{subsec:res-orient}).

The signal $s_i=\beta R(\psi_i)$ depends on its normal only through $\cos\psi_i=\vect n_i\trans\vect u$, so that the $i$th row of $\vect D$ is $\beta\dot R(\psi_i)\,\partial(\vect n_i\trans\vect u)/\partial\bm\eta\trans$. For independent errors, $\bm\eta=[\bm\eta_1\trans,\ldots,\bm\eta_K\trans]\trans$ and, to first order, $\vect n_i(\bm\eta)=\vect n_i+\vect E_i\bm\eta_i$, where the columns of $\vect E_i\in\R^{3\times2}$ form an orthonormal basis of the plane orthogonal to $\vect n_i$ and $\bm\eta_i\in\R^2$ holds the two tangent components of the error. The only nonzero entries of the $i$th row of $\vect D$ are then $\beta\dot R(\psi_i)\,\vect u\trans\vect E_i$ and, with $\bm\Sigma_\theta=\sigma_\theta^2\vect I_{2K}$, $\vect D\bm\Sigma_\theta\vect D\trans$ is diagonal with entries $\sigma_\theta^2\beta^2\dot R(\psi_i)^2\norm{\vect E_i\trans\vect u}^2$. Since $\vect E_i\vect E_i\trans=\vect I_3-\vect n_i\vect n_i\trans$, $\norm{\vect E_i\trans\vect u}^2=1-\cos^2\psi_i=\sin^2\psi_i$, which gives~\eqref{eq:effective-variance}. For a common error, $\bm\eta\in\R^3$ is a small rotation vector applied to all the normals, $\vect n_i(\bm\eta)=\vect n_i+\bm\eta\times\vect n_i$ to first order, which changes $\cos\psi_i$ by $(\bm\eta\times\vect n_i)\trans\vect u=\bm\eta\trans(\vect n_i\times\vect u)$. The rows of $\vect D$ are then $\beta\dot R(\psi_i)(\vect n_i\times\vect u)\trans$ and, with $\bm\Sigma_\theta=\sigma_\theta^2\vect I_3$, $\vect D\bm\Sigma_\theta\vect D\trans=\sigma_\theta^2\vect D\vect D\trans$ is a full matrix.
